% Topic 9: Matrices as Transformations
% Part of the Matrix Fundamentals section

\subsection{Matrices as Linear Transformations}

One of the most important applications of matrices is in representing geometric transformations. While you may be familiar with column vectors as representing points or translations, multiplying a vector by a matrix can scale, reflect, rotate, or shear that vector. 

A linear transformation can be thought of as a function that maps an input vector $\mathbf{v} = \begin{bmatrix} x \\ y \end{bmatrix}$ to an output vector $\mathbf{v'} = \begin{bmatrix} x' \\ y' \end{bmatrix}$ via matrix multiplication:
\[
\mathbf{v'} = A\mathbf{v} \quad \implies \quad \begin{bmatrix} x' \\ y' \end{bmatrix} = \begin{bmatrix} a & b \\ c & d \end{bmatrix} \begin{bmatrix} x \\ y \end{bmatrix}
\]

\subsubsection*{Common 2D Transformations}
Here are standard matrices for common geometric transformations in the 2D Cartesian plane:

\begin{center}
\renewcommand{\arraystretch}{1.5}
\begin{tabular}{p{2cm} p{4.5cm} p{6.5cm}}
\hline
\textbf{Transform} & \textbf{Matrix} & \textbf{Geometric Effect} \\
\hline
\textbf{Scaling} & 
$\begin{bmatrix} k_x & 0 \\ 0 & k_y \end{bmatrix}$ & 
Enlarges or shrinks shapes by factor $k_x$ in $x$-dir, $k_y$ in $y$-dir. \\

\textbf{Reflection} & 
$\begin{bmatrix} \cos 2\alpha & \sin 2\alpha \\ \sin 2\alpha & -\cos 2\alpha \end{bmatrix}$ & 
Reflects across a line through the origin making angle $\alpha$ with the positive $x$-axis ($y = x\tan\alpha$). \\

\textbf{Rotation} & 
$\begin{bmatrix} \cos\theta & -\sin\theta \\ \sin\theta & \cos\theta \end{bmatrix}$ & 
Rotates counter-clockwise by angle $\theta$ around the origin. \\

\textbf{Shear} & 
$\begin{bmatrix} 1 & k \\ 0 & 1 \end{bmatrix}$ (horizontal) \newline 
$\begin{bmatrix} 1 & 0 \\ k & 1 \end{bmatrix}$ (vertical) & 
Slants the shape along the axis by a factor $k$. \\
\hline
\end{tabular}
\end{center}

\subsubsection*{Composition of Transformations and Inverses}

If we apply a transformation $A$ and then a transformation $B$, this is equivalent to a single transformation given by the matrix product $BA$. Notice the order: $B(A\mathbf{v}) = (BA)\mathbf{v}$. 
Because matrix multiplication is generally \textbf{not commutative} ($AB \neq BA$), the order of transformations matters. For example, rotating an object and then reflecting it is usually different from reflecting it and then rotating it.

Furthermore, if a matrix $A$ represents a transformation, its inverse $A^{-1}$ (if it exists) represents \textbf{undoing} that transformation. The absolute value $|\det(A)|$ gives the area scale factor; the sign indicates orientation. If $\det(A) = 0$, the transformation collapses the 2D plane into a line or a point, making it non-invertible.

\begin{remark}[NSW HSC Heritage: Reflections and Rotations (1975--1980)]
A celebrated theorem examined in the NSW HSC 4~Unit examinations (notably 1975 Question~6 and 1976 Question~9) states that \textbf{the product of any two reflection matrices is a rotation matrix}.

Let $R_\alpha$ and $R_\beta$ be reflection matrices across lines inclined at angles $\alpha$ and $\beta$ respectively:
\[
R_\beta R_\alpha = \begin{bmatrix} \cos 2\beta & \sin 2\beta \\ \sin 2\beta & -\cos 2\beta \end{bmatrix} \begin{bmatrix} \cos 2\alpha & \sin 2\alpha \\ \sin 2\alpha & -\cos 2\alpha \end{bmatrix} = \begin{bmatrix} \cos 2(\beta - \alpha) & -\sin 2(\beta - \alpha) \\ \sin 2(\beta - \alpha) & \cos 2(\beta - \alpha) \end{bmatrix}.
\]
The composite transformation is a rotation through an angle $\theta = 2(\beta - \alpha)$, twice the angle between the two lines of reflection! Notice also that $\det(R_\alpha) = -\cos^2 2\alpha - \sin^2 2\alpha = -1$ (reversing orientation), while $\det(R_\beta R_\alpha) = (-1)(-1) = +1$ (preserving orientation), confirming geometrically why two reflections yield a rotation.
\end{remark}

\bigskip

\begin{problem}[Transformation Composition]
Let point $P$ have coordinates $(2, 3)$. 
\begin{enumerate}[label=\textbf{(\alph*)}]
    \item Find the coordinates of $P'$ after a counter-clockwise rotation of $90^\circ$ around the origin.
    \item Find the matrix that represents a reflection across the $x$-axis, followed by a counter-clockwise rotation of $90^\circ$.
    \item Use this combined matrix to transform point $P$.
\end{enumerate}
\end{problem}

\begin{solution}
\textbf{(a) Single Transformation} \\
The rotation matrix for $\theta = 90^\circ$ is $R_{90^\circ} = \begin{bmatrix} \cos 90^\circ & -\sin 90^\circ \\ \sin 90^\circ & \cos 90^\circ \end{bmatrix} = \begin{bmatrix} 0 & -1 \\ 1 & 0 \end{bmatrix}$.
Applying this to $P = \begin{bmatrix} 2 \\ 3 \end{bmatrix}$:
\[
P' = \begin{bmatrix} 0 & -1 \\ 1 & 0 \end{bmatrix} \begin{bmatrix} 2 \\ 3 \end{bmatrix} = \begin{bmatrix} -3 \\ 2 \end{bmatrix}
\]
So $P'$ is at $(-3, 2)$.

\medskip
\textbf{(b) Composition Matrix} \\
Let $M_1$ be the reflection across the $x$-axis: $M_1 = \begin{bmatrix} 1 & 0 \\ 0 & -1 \end{bmatrix}$. \\
Let $M_2$ be the $90^\circ$ rotation: $M_2 = \begin{bmatrix} 0 & -1 \\ 1 & 0 \end{bmatrix}$. \\
The combined transformation is $M = M_2 M_1$ (since $M_1$ acts first):
\[
M = \begin{bmatrix} 0 & -1 \\ 1 & 0 \end{bmatrix} \begin{bmatrix} 1 & 0 \\ 0 & -1 \end{bmatrix} = \begin{bmatrix} 0 & 1 \\ 1 & 0 \end{bmatrix}
\]
Notice that this resulting matrix corresponds to a reflection across the line $y = x$.

\medskip
\textbf{(c) Applying the Combined Matrix} \\
Applying $M$ to $P$:
\[
P'' = \begin{bmatrix} 0 & 1 \\ 1 & 0 \end{bmatrix} \begin{bmatrix} 2 \\ 3 \end{bmatrix} = \begin{bmatrix} 3 \\ 2 \end{bmatrix}
\]
So the final coordinate is $(3, 2)$.
\end{solution}

\begin{takeaways}
\item \textbf{Order Matters:} When chaining geometric transformations, you multiply the matrices from right to left.
\item \textbf{Geometric Meaning of Determinant:} If you apply the transformation matrix $M$ from part (b), $\det(M) = (0)(0) - (1)(1) = -1$. The absolute value of $1$ means the area is preserved, and the negative sign indicates that the orientation (e.g., clockwise vs counter-clockwise ordering of vertices) has been reversed, typical of a reflection.
\end{takeaways}
